%% kapa: exam=1 round=34 year=2023 period=2 form=A subject=appraisal_law subject_ko=감정평가관계법규 seq=21 number=21
%% source: Q-Net 감정평가사 기출문제 (cst00309) artlSeq=5211789 file=raw/1008/5211789/2023년도 제34회 감정평가사 1차시험 2교시.pdf
%% conversion: kordoc 4.24.51 -> tools/gen1.py (auto); review=proofread-sonnet; flags=none
\begin{question}{21}
국유재산법령상 국유재산에 관한 설명으로 옳은 것은?
\begin{choices}
  \item 재판상 화해에 의해 일반재산에 사권(私權)을 설정할 수 없다.
  \item 총괄청의 허가를 받은 경우라 할지라도 국유재산에 관한 사무에 종사하는 직원은 그 처리하는 국유재산을 취득할 수 없다.
  \item 국가 외의 자는 기부를 조건으로 축조하는 경우에도 국유재산에 영구시설물을 축조할 수 없다.
  \item 총괄청은 다음 연도의 국유재산의 관리·처분에 관한 계획의 작성을 위한 지침을 매년 6월 30일까지 중앙관서의 장에게 통보하여야 한다.
  \item 한국예탁결제원은 총괄청이나 중앙관서의 장등이 증권을 보관·취급하게 할 수 있는 법인에 해당한다.
\end{choices}
\end{question}
\answer{5}
